% Source: https://artofproblemsolving.com/community/c3h1567548_3d_geo_find_area
% Stable identity: aops-topic-1567548-q01
Use coordinates $A=(0,0,0)$, $B=(1,0,0)$, $C=(1,1,0)$, $D=(0,1,0)$ and $E,F,G,H$ one unit higher, so $G=(1,1,1)$. The inscribed sphere has centre $O=\left(\dfrac{1}{2},\dfrac{1}{2},\dfrac{1}{2}\right)$ and radius $\dfrac{1}{2}$.

\textbf{The plane.} The points $B$, $D$, $G$ all satisfy $x+y-z=1$, so this is the plane of triangle $BDG$, with normal vector $(1,1,-1)$.

\textbf{Distance from $O$.}
\[
\delta=\dfrac{\left|\dfrac{1}{2}+\dfrac{1}{2}-\dfrac{1}{2}-1\right|}{\sqrt{1+1+1}}=\dfrac{1}{2\sqrt3}.
\]
\textbf{The circle.} A plane at distance $\delta$ from the centre of a sphere of radius $R$ cuts it in a circle of radius $r$ with $r^2=R^2-\delta^2$ (Pythagorean theorem). Here
\[
r^2=\dfrac{1}{4}-\dfrac{1}{12}=\dfrac{1}{6},
\qquad\text{so the area is}\qquad \pi r^2=\boxed{\dfrac{\pi}{6}}.
\]

\begin{asy}
usepackage("amsmath");
size(6cm);
defaultpen(fontsize(10pt));
pair P(real x,real y,real z){return (x,z)+y*(0.5,0.36);}
pair A=P(0,0,0),B=P(1,0,0),C=P(1,1,0),D=P(0,1,0),Ep=P(0,0,1),F=P(1,0,1),G=P(1,1,1),H=P(0,1,1);
draw(A--B--F--Ep--cycle,linewidth(0.9));draw(B--C--G--F,linewidth(0.9));draw(Ep--H--G,linewidth(0.9));
draw(A--D,dashed);draw(D--C,dashed);draw(D--H,dashed);
fill(B--D--G--cycle,lightblue+opacity(0.5));
draw(B--D--G--cycle,blue+linewidth(0.9));
real k=1/sqrt(6);
real[] c={2/3,2/3,1/3};
real[] u={1/sqrt(2),-1/sqrt(2),0};real[] v={1/sqrt(6),1/sqrt(6),2/sqrt(6)};
path g;
for(int i=0;i<=72;++i){real t=2*pi*i/72;
  pair q=P(c[0]+k*(cos(t)*u[0]+sin(t)*v[0]),c[1]+k*(cos(t)*u[1]+sin(t)*v[1]),c[2]+k*(cos(t)*u[2]+sin(t)*v[2]));
  g=(i==0)?q:g--q;}
draw(g,red+linewidth(1));
dot(P(0.5,0.5,0.5));label("$O$",P(0.5,0.5,0.5),W);
dot(P(0.5,0.5,0),red);dot(P(1,0.5,0.5),red);dot(P(0.5,1,0.5),red);
label("$A$",A,SW);label("$B$",B,SE);label("$C$",C,E);label("$D$",D,W);label("$E$",Ep,NW);label("$F$",F,NW);label("$G$",G,NE);label("$H$",H,NW);
label("\mbox{}",(min(currentpicture,user=true)+max(currentpicture,user=true))/2,black); // keeps the text after the figure black when LaTeX typesets the labels inline
\end{asy}

\textbf{A check.} The sphere touches each face of the cube at the face's centre, which is the midpoint of that face's diagonals. So the sides $BD$, $DG$, $GB$ of triangle $BDG$ (face diagonals) touch the sphere at their midpoints (red dots), and the circle is the incircle of the equilateral triangle $BDG$ with side $\sqrt2$. Its inradius is $\dfrac{\sqrt2}{2\sqrt3}=\dfrac{1}{\sqrt6}$, giving the same area $\dfrac{\pi}{6}$.
